\section{习题}
\label{sec:eisenstein-exercises}
% ============================================================================

这一节的习题不是用来机械重复定义，而是帮助你把本章最重要的几条线索真正连起来。
前面我们先看见 Bernoulli 数如何出现在 $\zeta(k)$ 和 Eisenstein 级数的常数项里，接着用格求和构造 $G_k$，再借助 Poisson 求和把格点语言翻译成 $q$-展开，最后又把视野推广到 Dirichlet 特征与同余子群。
如果只读正文，这些步骤很容易看起来像一串彼此独立的技巧；习题的作用，正是让你亲手验证：它们其实是同一个故事的不同侧面。

做完这些题之后，你会更扎实地掌握三件事。
第一，你会真正会算 Bernoulli 数、$\zeta(2m)$ 和除数和，而不是只记住结论。
第二，你会对 Poisson 求和为什么能把解析信息变成算术信息有更直观的感觉。
第三，你会开始适应本章最重要的工作方式：先用对称性构造对象，再用 Fourier 展开读出系数，最后把这些系数解释成数论函数。
带着这个目标去做题，习题就不再是"附录"，而是理解 Eisenstein 级数最有效的第二遍学习。

\begin{exercise}
\label{ex:ch2-bernoulli-1}
证明 $B_1 = -1/2$，$B_2 = 1/6$，$B_4 = -1/30$。
（从递推关系 $\sum_{j=0}^{n-1}\binom{n}{j}B_j = 0$ 出发逐步计算。）
\end{exercise}

\begin{proof}[解答]
$n=2$：$\binom{2}{0}B_0 + \binom{2}{1}B_1 = 0 \Rightarrow 1 + 2B_1 = 0 \Rightarrow B_1 = -1/2$。

$n=3$：$B_0 + 3B_1 + 3B_2 = 0 \Rightarrow 1 - 3/2 + 3B_2 = 0 \Rightarrow B_2 = 1/6$。

$n=4$：$B_3 = 0$（$m$ 奇数 $\geq 3$ 时 $B_m = 0$），
$B_0 + 4B_1 + 6B_2 + 4B_3 + B_4 \cdot \frac{1}{5}\cdots$，
用 $n=5$：$1 + 5(-1/2) + 10(1/6) + 10 \cdot 0 + 5B_4 = 0$
$\Rightarrow 1 - 5/2 + 5/3 + 5B_4 = 0 \Rightarrow B_4 = -1/30$。
\end{proof}

\begin{exercise}
\label{ex:ch2-zeta-2}
用\cref{prop:bernoulli-zeta}计算 $\zeta(2), \zeta(4), \zeta(6)$，
验证 $\zeta(2) = \pi^2/6$。
\end{exercise}

\begin{proof}[解答]
$\zeta(2) = -(2\pi i)^2 B_2 / (2 \cdot 2!) = -(-4\pi^2)(1/6)/4 = \pi^2/6$ \checkmark。

$\zeta(4) = -(2\pi i)^4 B_4/(2 \cdot 24) = -(16\pi^4)(-1/30)/48 = \pi^4/90$ \checkmark。

$\zeta(6) = -(2\pi i)^6 B_6/(2 \cdot 720) = -(-64\pi^6)(1/42)/1440 = \pi^6/945$ \checkmark。
\end{proof}

\begin{exercise}
\label{ex:ch2-sigma-multiplicative}
证明 $\sigma_{k-1}(n)$ 是\textbf{积性函数}：$\gcd(m,n)=1 \Rightarrow \sigma_{k-1}(mn) = \sigma_{k-1}(m)\sigma_{k-1}(n)$。

（提示：$mn$ 的正因子恰好是 $m$ 的因子与 $n$ 的因子的乘积。）
\end{exercise}

\begin{proof}[解答]
若 $\gcd(m,n)=1$，则 $mn$ 的每个因子 $d$ 唯一分解为 $d = d_1 d_2$，其中 $d_1 \mid m$，$d_2 \mid n$，$\gcd(d_1, d_2) = 1$。故
\[
  \sigma_{k-1}(mn) = \sum_{d \mid mn} d^{k-1}
  = \sum_{d_1 \mid m}\sum_{d_2 \mid n} (d_1 d_2)^{k-1}
  = \left(\sum_{d_1 \mid m} d_1^{k-1}\right)\left(\sum_{d_2 \mid n} d_2^{k-1}\right)
  = \sigma_{k-1}(m)\sigma_{k-1}(n). \qedhere
\]
\end{proof}

\begin{exercise}
\label{ex:ch2-poisson-gaussian}
设 $f_t(x) = e^{-\pi t x^2}$（$t > 0$）。
\begin{enumerate}
  \item 证明 $\hat{f}_t(\xi) = t^{-1/2} e^{-\pi\xi^2/t}$。
  \item 用 Poisson 求和推导 $\theta(t) = t^{-1/2}\theta(1/t)$，
    其中 $\theta(t) = \sum_{n \in \Z} e^{-\pi n^2 t}$。
\end{enumerate}
\end{exercise}

\begin{proof}[解答]
(1) $\hat{f}_t(\xi) = \int e^{-\pi tx^2 - 2\pi i\xi x}\,dx$。配方：
$-\pi t x^2 - 2\pi i\xi x = -\pi t(x + i\xi/t)^2 - \pi\xi^2/t$。
故 $\hat{f}_t(\xi) = e^{-\pi\xi^2/t} \int e^{-\pi t(x+i\xi/t)^2}\,dx = t^{-1/2} e^{-\pi\xi^2/t}$
（围道积分，同\cref{ex:gaussian-fourier}）。

(2) $\theta(t) = \sum_{n \in \Z} f_t(n)$。
由 Poisson 求和（\cref{thm:poisson-summation}），
$\sum_n f_t(n) = \sum_n \hat{f}_t(n) = t^{-1/2}\sum_n e^{-\pi n^2/t} = t^{-1/2}\theta(1/t)$。
\end{proof}

\begin{exercise}
\label{ex:ch2-dirichlet-L-s1}
设 $\chi$ 是模 $N$ 的非主特征。证明 $L(1,\chi) \neq 0$。

（提示：先证 $\prod_\chi L(s,\chi) = \prod_p \prod_\chi (1 - \chi(p)p^{-s})^{-1}$
当 $s \to 1^+$ 时趋向无穷大；若某个 $L(1,\chi_0^{}) = 0$ 则乘积有界，矛盾。）
\end{exercise}

\begin{proof}[解答（思路）]
令 $s > 1$ 实数。考虑乘积
$\prod_{\chi \bmod N} L(s, \chi)$，其中 $\chi$ 跑遍所有模 $N$ 特征。
利用 Euler 乘积与特征的正交性，可以证明
\[
  \prod_{\chi} L(s, \chi) = \prod_{p} \frac{1}{\prod_\chi(1 - \chi(p)p^{-s})}.
\]
用 $\sum_\chi \chi(p^k) \geq 0$（群特征的正半定性），
整个乘积 $\geq 1$ 且当 $s \to 1^+$ 时由 $L(s, \chi_0) \sim \frac{\varphi(N)}{N(s-1)}$（极点）趋向无穷。

若 $L(1, \chi_1) = 0$ 对某非主 $\chi_1$，则 $L(1, \bar\chi_1) = \overline{L(1,\chi_1)} = 0$，
两个零点与主特征的极点"对消"，导致整个乘积在 $s = 1$ 处有界，矛盾。
详细证明见 Serre《算术导引》或 Davenport《乘性数论》。
\end{proof}
